\title{DeepSeekMath: Pushing the Limits of Mathematical Reasoning in Open Language Models}

\begin{abstract}
Language models often struggle with mathematical reasoning due to the inherently complex and highly structured characteristics of mathematical content. In this study, we present DeepSeekMath~7B, which extends the pre-training of DeepSeek-Coder-Base-v1.5 7B by incorporating 120B math-focused tokens, compiled from Common Crawl as well as natural language and coding sources. DeepSeekMath~7B attains a strong result of 51.7\% on the MATH benchmark at the competition level, all without dependence on external toolkits or voting-based strategies, thereby nearing the effectiveness demonstrated by Gemini-Ultra and GPT-4. When employing self-consistency over 64 generations, DeepSeekMath~7B reaches a performance of 60.9\% on MATH. The advancement in mathematical reasoning delivered by DeepSeekMath~relies on two principal innovations: firstly, the model leverages abundant public web data via a carefully crafted data selection process; secondly, we propose Group Relative Policy Optimization (GRPO), a modification of Proximal Policy Optimization (PPO), which not only strengthens mathematical reasoning skills but also significantly improves PPO’s memory efficiency.
\end{abstract}

\section{Introduction}

Large language models (LLM) have transformed methods of mathematical reasoning within artificial intelligence, driving notable progress across both quantitative reasoning benchmarks \citep{MATH} and geometry reasoning benchmarks \citep{trinh2024solving}. In addition, these models serve as valuable tools in supporting humans as they tackle challenging mathematical problems \citep{tao}. Despite these advancements, state-of-the-art models like GPT-4 \citep{gpt4} and Gemini-Ultra \citep{gemini} remain inaccessible to the public, and the best available open-source alternatives exhibit significantly lower performance.

In this paper, we present DeepSeekMath, a language model tailored to the mathematical domain that markedly surpasses other open-source alternatives in mathematical problem-solving and closely matches the performance of GPT-4 on academic tests. Our approach involves constructing the DeepSeekMath Corpus, an extensive and high-quality pre-training dataset consisting of 120 billion tokens focused on mathematics. The data is sourced from the Common Crawl (CC) and filtered through a fastText-based classifier \citep{joulin2016fasttext}. Initially, the classifier is developed using examples from OpenWebMath \citep{openwebmath} as the positive class, supplemented with a broad set of unrelated web content as negatives. This classifier is then applied to CC to identify further positive samples, which are subsequently curated with human annotation to improve reliability. The classifier undergoes a retraining step using the expanded and refined dataset to boost classification quality. Evaluations demonstrate that the resultant corpus maintains exceptional quality: our foundational model, DeepSeekMath-Base 7B, attains 64.2\% on the GSM8K benchmark \citep{gsm8k} and 36.2\% on the advanced MATH benchmark \citep{MATH}, thereby exceeding the results obtained by Minerva 540B \citep{minerva}. Furthermore, as the DeepSeekMath Corpus supports multiple languages, we observe enhanced outcomes on Chinese mathematical benchmarks \citep{wei2023cmath,agieval}. We regard our data-processing methodology as an initial contribution for future exploration within the community and anticipate considerable potential for further advances.

DeepSeekMath-Base utilizes DeepSeek-Coder-Base-v1.5 7B \citep{deepseek-coder} as its initialization point, based on our observation that models pre-trained on code tend to yield better results than those derived from general LLMs. Additionally, we find that mathematical training enhances not only performance on mathematical tasks but also leads to improvements on the MMLU \citep{mmlu} and BBH \citep{bbh} benchmarks. This suggests that mathematical instruction benefits broader reasoning abilities beyond mathematics.

Following the pre-training stage, we fine-tune DeepSeekMath-Base using mathematical instruction datasets that include chain-of-thought \citep{cot}, program-of-thought \citep{pot,pal}, and tool-integrated reasoning \citep{tora} examples. The final model, DeepSeekMath-Instruct 7B, surpasses all other 7B-sized models and performs on par with open-source instruction-tuned models of 70B scale.

In addition, we propose Group Relative Policy Optimization (GRPO), an alternative reinforcement learning (RL) approach inspired by Proximal Policy Optimization (PPO) \citep{schulman2017proximal}. Unlike traditional PPO, GRPO eliminates the need for a critic model and instead derives the baseline directly from group scores. This modification leads to a marked reduction in the computational cost associated with training. Utilizing only a portion of English instruction tuning datasets, GRPO delivers pronounced gains over the robust DeepSeekMath-Instruct model, enhancing performance in both in-domain evaluations (GSM8K: 82.9\% $\rightarrow$ 88.2\%, MATH: 46.8\% $\rightarrow$ 51.7\%) as well as on out-of-domain mathematical benchmarks (e.g., CMATH: 84.6\% $\rightarrow$ 88.8\%) in the RL stage. We further introduce a comprehensive framework that unifies various methodologies, including Rejection Sampling Fine-Tuning (RFT) \citep{yuan2023scaling}, Direct Preference Optimization (DPO) \citep{dpo}, PPO, and GRPO, showing that they can all be interpreted as either explicit or streamlined RL approaches. Through a broad array of experimental analyses—such as contrasting online versus offline training, comparing outcome against process supervision, and evaluating single-turn with iterative RL—we systematically explore the fundamental attributes of this shared framework. Finally, we clarify the mechanisms by which our RL techniques enhance instruction-tuned models and outline several prospective research avenues for further optimizing RL efficacy within this unified perspective.

\subsection{Contributions}
We present advances in scalable math pre-training, complemented by an in-depth investigation and evaluation of reinforcement learning methodologies.

\textbf{Math Pre-Training at Scale}

\begin{itemize}[topsep=0pt]
    \item In this study, we demonstrate that the openly available Common Crawl dataset holds significant mathematical content. Utilizing a carefully crafted data selection process, we assembled the DeepSeekMath Corpus—a comprehensive dataset comprising 120B tokens sourced from web pages specifically curated for mathematical relevance. This collection surpasses the scale of the math-related web pages employed by Minerva \citep{minerva} by nearly sevenfold, and is approximately nine times larger than the newly introduced OpenWebMath dataset \citep{openwebmath}.

\item DeepSeekMath-Base 7B, our foundational model trained in advance, demonstrates performance on par with Minerva 540B \citep{minerva}. This suggests that mathematical reasoning ability is not solely determined by model size. Robust results can also be attained by pre-training smaller models on high-quality datasets.

\item We present results from our experiments involving math training. Pre-training models on code enhances their capacity to tackle mathematical tasks, regardless of whether tools are used. This provides some evidence regarding the enduring question: \textit{does code training enhance reasoning skills?} Our results suggest that, specifically for mathematical reasoning, it indeed has a positive effect.

\item While it is standard practice to utilize arXiv papers for training—particularly within numerous mathematics-focused works—doing so does not yield significant gains across any of the mathematical benchmarks used in this study.
\end{itemize}

\textbf{Exploration and Analysis of Reinforcement Learning}

\begin{itemize}[topsep=0pt]
    \item We propose Group Relative Policy Optimization (GRPO), a reinforcement learning algorithm designed for both efficiency and effectiveness. Unlike traditional methods that rely on a critic model, GRPO leverages group scores to estimate the baseline, leading to substantial savings in computational resources over Proximal Policy Optimization (PPO).
    \item Our results indicate that employing GRPO yields substantial improvements in the instruction-tuned model DeepSeekMath-Instruct when trained solely on instruction-tuning data. Additionally, we find that reinforcement learning with GRPO results in better generalization to out-of-domain tasks.
    \item We establish a unified framework to interpret various algorithms, including RFT, DPO, PPO, and GRPO. Comprehensive experimental analyses are carried out, such as comparisons between online and offline training paradigms, supervision via outcomes versus process, and the effects of single-turn versus iterative reinforcement learning, in order to rigorously examine the key components within this framework.
    \item Drawing from this unified framework, we investigate the underlying factors that drive the success of reinforcement learning and identify promising avenues for further improving the effectiveness of reinforcement learning with LLMs.
\end{itemize}

\subsection{Summary of Evaluations and Metrics}

\begin{itemize}[topsep=0pt]
    \item \textbf{English and Chinese Mathematical Reasoning}: Our evaluation involves an extensive suite of English and Chinese benchmarks, encompassing mathematical tasks that span from elementary to university curricula. The English datasets assessed consist of GSM8K \citep{gsm8k}, MATH \citep{MATH}, SAT \citep{llemma}, OCW Courses \citep{minerva}, and MMLU-STEM \citep{mmlu}. For the Chinese language, we utilize benchmarks such as MGSM-zh \citep{mgsm}, CMATH \citep{wei2023cmath}, Gaokao-MathCloze \citep{agieval}, and Gaokao-MathQA \citep{agieval}. Our evaluation framework measures model proficiency in producing complete textual solutions unaided by external tools, alongside their capacity to tackle problems through Python-based computations.

On English evaluation metrics, DeepSeekMath-Base demonstrates performance on par with the proprietary Minerva 540B model \citep{minerva}, and outperforms every available open-source base model—including Mistral 7B \citep{mistral} and Llemma-34B \citep{llemma}—whether or not they have been subjected to math-specific pre-training, often by a wide margin. Importantly, DeepSeekMath-Base achieves superior results on Chinese evaluation sets, which can be attributed to our departure from the English-only data collection strategies seen in previous works \citep{minerva, llemma}. Instead, our approach incorporates high-quality mathematical data from multiple languages. After applying both mathematical instruction tuning and reinforcement learning, the advanced versions, DeepSeekMath-Instruct and DeepSeekMath-RL, deliver robust results—exceeding 50\% accuracy on the competitive MATH dataset, marking the first time this has been achieved by an open-source model.

\item \textbf{Formal Mathematics}: DeepSeekMath-Base is assessed through the informal-to-formal theorem proving benchmark introduced by \citep{dsp_proof}, utilizing miniF2F \citep{minif2f} as the dataset and Isabelle \citep{isabelle} as the selected proof assistant. The model attains notable results in few-shot autoformalization settings.
\item \textbf{Natural Language Understanding, Reasoning, and Code}: To thoroughly examine the general capabilities of models in natural language comprehension, reasoning, and programming, DeepSeekMath-Base is evaluated using the Massive Multitask Language Understanding (MMLU) benchmark \citep{mmlu}, featuring 57 diverse multiple-choice subjects, as well as BIG-Bench Hard (BBH) \citep{bbh}, a collection of 23 advanced tasks that demand complex, multi-step reasoning. In addition, we include HumanEval \citep{codex} and MBPP \citep{mbpp} to assess programming aptitude, as both are established code evaluation benchmarks. The math-centric pre-training notably enhances outcomes in both understanding and reasoning domains.

\section{Math Pre-Training}

\subsection{Data Collection and Decontamination} Here, we describe the methodology used to assemble the DeepSeekMath Corpus utilizing data from Common Crawl. Figure \ref{fig:data_collect} illustrates the iterative pipeline developed for the systematic extraction of a comprehensive mathematical dataset, beginning with an initial, high-quality seed collection containing math-specific content. Importantly, this framework is not restricted to mathematics; it is equally adaptable for building large-scale datasets in other specialized fields, such as programming.

We begin by selecting OpenWebMath \citep{openwebmath}, a curated repository of high-quality mathematical documents from the web, to serve as our initial seed dataset. Utilizing this seed set, we train a fastText model \citep{joulin2016fasttext} with the goal of retrieving additional web pages similar in style and content to OpenWebMath. For the training process, 500,000 samples are randomly drawn from the seed corpus as positive instances, while another 500,000 web pages are randomly sampled from Common Crawl and treated as negative instances. Model training is conducted using a publicly available library\footnote{\url{https://fasttext.cc}}, with the following hyperparameters: a vector size of 256, a learning rate set at 0.1, a maximum word n-gram length of 3, a minimum frequency threshold of 3, and 3 epochs for training. In order to manage the large scale of Common Crawl, we apply deduplication techniques at the URL level as well as near-duplicate content filtering, which reduces the corpus to 40 billion HTML documents. The trained fastText model is then used to retrieve web pages with mathematical content from this deduplicated dataset. To filter for higher-quality mathematics content, we rank candidate pages based on their fastText model scores and retain only those with the highest scores. The size of the data selected is further evaluated via pre-training trials on subsets containing the top 40B, 80B, 120B, and 160B tokens. In this initial round, the top 40B tokens are preserved.

Following the initial round of data gathering, a substantial number of mathematical web pages remain uncollected, largely due to the limited diversity of the positive examples used to train the fastText model. To enhance the model's performance, we expand the seed set by sourcing additional mathematical web domains. In particular, the full Common Crawl corpus is partitioned into non-overlapping domains, where each domain consists of web pages with an identical base URL. For every domain, we determine the fraction of web pages collected in the first pass. Domains with more than 10\% of their pages captured during the first iteration are designated as math-related (for example, \url{mathoverflow.net}). Next, we conduct manual annotation to label URLs within these domains that contain mathematical content (e.g., \url{mathoverflow.net/questions}). Pages connected to these URLs that were not initially collected are incorporated into the expanded seed corpus. This strategy increases the pool of positive samples, enabling the fastText model to better retrieve mathematical content in subsequent rounds. After four cycles of data acquisition, our process yields a total of 35.5M mathematical pages, amassing 120B tokens. Upon reaching the fourth iteration, we observe that almost 98\% of relevant data had already been obtained during the third cycle, leading us to halt further collection.

In order to prevent benchmark contamination, we adopt the approach described in \cite{deepseek-coder} to exclude web pages that feature questions or solutions derived from established English mathematical benchmarks like GSM8K \citep{gsm8k} and MATH \citep{MATH}, as well as Chinese resources such as CMATH \citep{wei2023cmath} and AGIEval \citep{agieval}. Specifically, we eliminate any portion of text within our mathematical training dataset that contains a sequence of 10 consecutive words precisely matching any substring from these evaluation benchmarks. For benchmark passages whose lengths are below 10 grams but at least 3 grams, we apply an exact-match filtering method to ensure contaminated webpages are also discarded.

\subsection{Validating the Quality of the DeepSeekMath~Corpus}

We run pre-training experiments to investigate how the DeepSeekMath~Corpus is compared with the recently released math-training corpora:

\begin{itemize}[topsep=0pt]
    \item \textbf{MathPile} \citep{mathpile}: This dataset compiles material from a range of sources—textbooks, Wikipedia, ProofWiki, CommonCrawl, StackExchange, and arXiv—accumulating to 8.9B tokens, with arXiv papers making up more than 85\% of the total corpus;
    \item \textbf{OpenWebMath} \citep{openwebmath}: Constructed by extracting mathematical texts from the CommonCrawl archive, resulting in a collection with 13.6B tokens;
    \item \textbf{Proof-Pile-2} \citep{llemma}: This composite mathematical dataset comprises OpenWebMath, AlgebraicStack (providing 10.3B tokens of mathematical code), and arXiv articles (amounting to 28.0B tokens). In alignment with \cite{llemma}, our experiments utilizing Proof-Pile-2 maintain an arXiv:Web:Code distribution ratio of 2:4:1.
\end{itemize}

\subsubsection{Training Setting}
\label{sec:quality-policy}
We fine-tune a general-purpose language model with 1.3B parameters—constructed on the DeepSeek LLM architecture \citep{deepseek-llm} and identified as DeepSeek-LLM 1.3B—on various mathematical datasets individually, each for a total of 150B tokens. The training is performed with the streamlined and resource-efficient HAI-LLM training suite \citep{haillm}. Adhering to the optimization protocol established for DeepSeek LLMs, we make use of the AdamW optimizer \citep{adamW} configured with $\beta_1=0.9$, $\beta_2=0.95$, and $\mathrm{weight\_decay}=0.1$. The learning rate is governed by a multi-phase schedule: after a warmup phase lasting 2,000 steps, it ascends to its maximum, then declines to 31.6\% of its peak at the 80\% mark of training, and further drops to 10.0\% of the peak by the end of 90\% of the process. The highest learning rate adopted during training is 5.3e-4. We utilize a batch size consisting of 4M tokens, paired with a context window of 4K tokens.

\subsubsection{Evaluation Results}

\textbf{The DeepSeekMath~Corpus is of high quality, covers multilingual mathematical content, and is the largest in size.}

\begin{itemize}[topsep=0pt]
    \item \textbf{High-quality}: 
    Downstream efficacy is assessed across eight mathematical evaluation tasks using few-shot chain-of-thought prompting, as detailed in \cite{cot}. Table \ref{tab:corpora_comparison} illustrates that models trained with the DeepSeekMath~Corpus outperform others by a notable margin. Furthermore, Figure \ref{fig:corpora_comparisons} reveals that at 50B tokens (corresponding to a complete epoch for Proof-Pile-2), the DeepSeekMath~Corpus-based model surpasses Proof-Pile-2, underscoring its superior average data quality.
    \item \textbf{Multilingual}:
    DeepSeekMath~Corpus contains multilingual datasets, with English and Chinese serving as the principal languages. Table \ref{tab:corpora_comparison} shows that training on the DeepSeekMath~Corpus leads to improved mathematical reasoning capabilities in both English and Chinese, whereas previously existing mathematical corpora, largely focused on English, either provide minimal benefit or may even detract from Chinese mathematical reasoning performance.
    \item \textbf{Large-scale}:
    With a significantly greater volume of content compared to prior mathematical corpora, the DeepSeekMath~Corpus offers an order-of-magnitude increase in dataset size. As depicted in Figure \ref{fig:corpora_comparisons}, DeepSeek-LLM 1.3B trained on DeepSeekMath~Corpus demonstrates both a sharper trajectory in learning progress and more sustained performance improvements. Conversely, the smaller baseline corpora—having undergone multiple epochs—lead to rapid stagnation in model capabilities.
\end{itemize}

\subsection{Training and Evaluating DeepSeekMath-Base 7B}

This section details the development of DeepSeekMath-Base 7B, a foundational model designed with advanced reasoning capabilities, particularly suited to mathematical tasks. The model’s parameters are initialized using DeepSeek-Coder-Base-v1.5 7B \citep{deepseek-coder}, and it undergoes training on a corpus of 500B tokens. The dataset comprises 56\% DeepSeekMath Corpus, 4\% AlgebraicStack, 10\% arXiv material, 20\% code sourced from Github, and 10\% natural language data drawn from English and Chinese Common Crawl. Training largely follows the approach described in Section \ref{sec:quality-policy}; however, the peak learning rate is set at 4.2e-4 and each batch consists of 10M tokens.

We present an extensive evaluation of DeepSeekMath-Base 7B's proficiency in mathematics, analyzing its capacity to independently generate complete mathematical solutions, utilize external tools for problem-solving, and engage in formal theorem proving. In addition to mathematical competencies, we also examine the model’s broader characteristics, highlighting its abilities in natural language comprehension, reasoning, and programming.

\paragraph{Mathematical Problem Solving with Step-by-Step Reasoning}

We assess the ability of DeepSeekMath-Base to solve mathematical problems using few-shot chain-of-thought prompting \citep{cot} across eight different benchmarks in both English and Chinese. These benchmarks include a variety of quantitative reasoning datasets (such as GSM8K \citep{gsm8k}, MATH \citep{MATH}, and CMATH \citep{wei2023cmath}), as well as multiple-choice assessments (like MMLU-STEM \citep{mmlu} and Gaokao-MathQA \citep{agieval}), spanning a wide array of mathematical disciplines from elementary through to university-level tasks.

Table \ref{tab:base_math_cot} demonstrates that DeepSeekMath-Base 7B outperforms all other open-source base models across the eight evaluated benchmarks. This comparison includes notable models such as the broadly adopted general model Mistral 7B \citep{mistral} and the recently introduced Llemma 34B \citep{llemma}, which received mathematics-specific training on Proof-Pile-2 \citep{llemma}. Particularly on the MATH dataset, which assesses competition-level mathematics, DeepSeekMath-Base achieves a more than 10\% absolute improvement over any prior open-source base model, and even exceeds the performance of Minerva 540B \citep{minerva}, a closed-source base model built atop PaLM \citep{palm}, which is 77 times larger and also specialized on mathematical corpora.

\paragraph{Mathematical Problem Solving with Tool Use} We assess the ability of models to solve mathematical problems by leveraging tools on the GSM8K and MATH benchmarks, employing a few-shot program-of-thought prompting approach \citep{pot,pal}. For each task, the model is instructed to generate a Python script, making use of libraries like \textit{math} and \textit{sympy} to perform complex calculations as needed. The outcome of running the generated program is used as the solution. As detailed in Table \ref{tab:base_math_pot_proof}, DeepSeekMath-Base 7B achieves superior results compared to the previously leading Llemma 34B model.

\paragraph{Formal Mathematics}

Automating the process of formal proof construction plays a pivotal role in verifying the correctness of mathematical arguments and improving productivity, a topic that has garnered significant scholarly focus in recent times. In this work, we assess the capabilities of DeepSeekMath-Base 7B in the informal-to-formal proof generation task presented by \citep{dsp_proof}, which involves producing a formal proof given an informal problem description, its formal expression, and an informal argument. Our evaluation uses the miniF2F dataset \citep{minif2f}, a recognized benchmark designed to test formalization in Olympiad-level mathematics. For each item, we employ few-shot prompting to elicit Isabelle-formalized proofs. Consistent with the methodology of \cite{dsp_proof}, our approach entails producing high-level proof outlines with the model and subsequently applying Sledgehammer \citep{sledgehammer}, a standard automated proving tool, to elaborate any omitted steps. As summarized in Table \ref{tab:base_math_pot_proof}, DeepSeekMath-Base 7B achieves notable results in automatic proof formalization.

\paragraph{Natural Language Understanding, Reasoning, and Code}

We assess the capabilities of the model in natural language understanding using MMLU \citep{mmlu}, evaluate reasoning on the BBH dataset \citep{bbh}, and test programming proficiency with HumanEval \citep{codex} and MBPP \citep{mbpp}. Table \ref{tab:understanding_reasoning_code} demonstrates that DeepSeekMath-Base 7B delivers marked improvements in both MMLU and BBH scores when compared to its predecessor, DeepSeek-Coder-Base-v1.5 \citep{deepseek-coder}, thus demonstrating the beneficial effects of mathematical training on understanding and reasoning tasks. Furthermore, the inclusion of code tokens in continued training allows DeepSeekMath-Base 7B to largely preserve the coding performance levels observed in DeepSeek-Coder-Base-v1.5 on both coding evaluation sets. Taken together, these results indicate that DeepSeekMath-Base 7B surpasses the general-purpose model Mistral 7B \citep{mistral} across all three evaluated benchmarks involving reasoning and code generation.

\section{Supervised Fine-Tuning}

\subsection{SFT Data Curation}

We develop a dataset for mathematical instruction-tuning that encompasses both English and Chinese problems drawn from a variety of mathematical domains and featuring a range of difficulty levels. Each problem is accompanied by a corresponding solution presented in chain-of-thought (CoT) \citep{cot}, program-of-thought (PoT) \citep{pot,pal}, or tool-integrated reasoning format \citep{tora}. Altogether, the dataset comprises 776K training instances.

\begin{itemize}[topsep=0pt]
    \item \textbf{English mathematical datasets}: We provide tool-based solution annotations for problems from GSM8K and MATH, and utilize a subset of MathInstruct \citep{MathInstruct} together with the training split of Lila-OOD \citep{lila}, where answers are derived via CoT or PoT methods. The English dataset spans multiple mathematical domains, including algebra, probability, number theory, calculus, and geometry.
    \item \textbf{Chinese mathematical datasets}: Our collection comprises K-12 mathematics questions in Chinese that cover 76 different sub-topics (e.g., linear equations), each with solutions in both CoT and tool-integrated reasoning formats.
\end{itemize}

\subsection{Training and Evaluating DeepSeekMath-Instruct 7B}

This section presents DeepSeekMath-Instruct 7B, a model refined from DeepSeekMath-Base through mathematical instruction tuning. For training, we randomly concatenate training samples to fill a context window up to 4K tokens. The model is optimized over 500 steps, using a batch size of 256 and maintaining a fixed learning rate of 5e-5 throughout the process.

We evaluate models' mathematical performance both without and with tool use, on 4 quantitative reasoning benchmarks in English and Chinese. We benchmark our model against the leading models of the time:

\begin{itemize}[topsep=0pt]
    \item \textbf{Closed-source models} encompass: (1) the GPT series, notably GPT-4 \citep{gpt4} and the GPT-4 Code Interpreter~\footnote{\url{https://openai.com/blog/chatgpt-plugins\#\#code-interpreter}} recognized as their most advanced variants, (2) Gemini Ultra and Pro \citep{gemini}, (3) Inflection-2 \citep{inflection-2}, and (4) Grok-1~\footnote{\url{https://x.ai/model-card}}, as well as several models from Chinese enterprises, such as (5) Baichuan-3~\footnote{\url{https://www.baichuan-ai.com}} and (6) the recent GLM-4~\footnote{\url{https://open.bigmodel.cn/dev/api\#glm-4}}

    from the GLM family \citep{glm}.

These models are designed for a wide array of applications, and the majority have been subjected to various alignment processes.

\item \textbf{Open-source models} in this context encompass: foundational models such as (1) DeepSeek-LLM-Chat 67B \citep{deepseek-llm}, (2) Qwen 72B \citep{qwen}, (3) SeaLLM-v2 7B \citep{seallm}, and (4) ChatGLM3 6B \citep{chatglm3}; in addition, models with specialized mathematical capabilities are represented by (5) InternLM2-Math 20B~\footnote{\url{https://github.com/InternLM/InternLM-Math}}, which is based on InternLM2 and further trained on mathematical corpora followed by instruction fine-tuning, (6) Math-Shepherd-Mistral 7B, which applies PPO training \citep{schulman2017proximal} to Mistral 7B \citep{mistral} utilizing a process-supervised reward framework, (7) the WizardMath family \citep{wizardmath}, which enhances Mistral 7B and Llama-2 70B \citep{llama2} for mathematical reasoning by means of evolve-instruct (a variant of instruction tuning leveraging AI-generated instructions) and PPO optimization, predominantly using data from GSM8K and MATH, (8) MetaMath 70B \citep{metamath}, a version of Llama-2 70B further trained on an expanded set of problems from GSM8K and MATH, (9) ToRA 34B \cite{tora}, which adapts CodeLlama 34B for tool-assisted mathematical reasoning via fine-tuning, and (10) MAmmoTH 70B \citep{MathInstruct}, which results from instruction-tuning Llama-2 70B on MathInstruct data.

Table \ref{tab:sft_rl_math} illustrates that, when evaluated without the allowance of external tool use, DeepSeekMath-Instruct 7B achieves notable proficiency in stepwise mathematical reasoning. Remarkably, on the challenging competition-level MATH benchmark, our model outperforms all open-source alternatives as well as most commercial systems (such as Inflection-2 and Gemini Pro) by no less than 9\% in absolute accuracy. This advantage is maintained even in comparison to models with far larger parameter counts (for example, Qwen 72B) and those specifically optimized through math-centered reinforcement learning, like WizardMath-v1.1 7B. While DeepSeekMath-Instruct matches the performance of leading Chinese proprietary systems GLM-4 and Baichuan-3 on the MATH dataset, it continues to lag behind the leading models GPT-4 and Gemini Ultra.

In the evaluation context that permits the combination of natural language reasoning with programmatic tool application for problem solving, DeepSeekMath-Instruct 7B achieves nearly 60\% accuracy on the MATH benchmark, outperforming all currently available open-source models. Across additional benchmarks, our model demonstrates comparable performance to DeepSeek-LLM-Chat 67B, the previous leading model, despite being only one-tenth its size.

\subsection{Group Relative Policy Optimization} Reinforcement learning (RL) has demonstrated considerable success in enhancing the mathematical reasoning skills of LLMs following the Supervised Fine-Tuning (SFT) phase \citep{wang2023math,wizardmath}. Here, we present Group Relative Policy Optimization (GRPO), our proposed RL approach designed to be both efficient and effective.

\subsubsection{From PPO to GRPO} Proximal Policy Optimization (PPO) \citep{schulman2017proximal} is an actor-critic RL algorithm that is widely used in the RL fine-tuning stage of LLMs \citep{ouyang2022training}. In particular, it optimizes LLMs by maximizing the following surrogate objective:

\begin{equation}
\footnotesize
    \mathcal{J}_{PPO}(\theta) = \mathbb{E}{[q \sim P(Q), o \sim \pi_{\theta_{old}}(O|q)]} \frac{1}{|o|} \sum_{t=1}^{|o|} \min \left[ \frac{\pi_\theta(o_{t} | q, o_{<t})}{\pi_{\theta_{old}}(o_{t} | q, o_{<t})} A_{t}, \text{clip} \left( \frac{\pi_\theta(o_{t} | q, o_{<t})}{\pi_{\theta_{old}}(o_{t} | q, o_{<t})}, 1 - \varepsilon, 1 + \varepsilon \right)  A_{t} \right] ,
\end{equation}

where $\pi_{\theta}$ and $\pi_{\theta_{old}}$ denote the current and previous policy models, while $q$ and $o$ represent questions and corresponding outputs that are sampled from both the question dataset and the prior policy $\pi_{\theta_{old}}$. The term $\varepsilon$ refers to a hyper-parameter associated with clipping in PPO, which serves to promote the stability of the training process. $A_t$ stands for the advantage, obtained using Generalized Advantage Estimation (GAE) \citep{gae}, leveraging the future rewards $\{r_{\ge t}\}$ and a value function $V_{\psi}$ that is learned during training. Consequently, PPO requires simultaneous training of a value function in addition to the policy network. To address potential issues arising from excessive reward model optimization, a typical solution is to incorporate a per-token KL divergence penalty against a reference model into the token-level reward, as described in \citep{ouyang2022training}, specifically:

\begin{equation}
    r_{t} = r_\varphi(q, o_{\le t}) - \beta \log\frac{\pi_{\theta}(o_{t}|q, o_{<t})}{\pi_{ref}(o_{t}|q, o_{<t})},
\label{eq:PPO-reward}
\end{equation}

Here, $r_\varphi$ denotes the reward model, $\pi_{ref}$ refers to the reference model—typically the initial SFT model—and $\beta$ represents the scaling factor applied to the KL divergence penalty.

As the value function employed in PPO is typically another model of comparable size as the policy model, it brings a substantial memory and computational burden. Additionally, during RL training, the value function is treated as a baseline in the calculation of the advantage for variance reduction. While in the LLM context, usually only the last token is assigned a reward score by the reward model, which may complicate the training of a value function that is accurate at each token. To address this, as shown in Figure \ref{fig:grpo}, we propose Group Relative Policy Optimization (GRPO), which obviates the need for additional value function approximation as in PPO, and instead uses the average reward of multiple sampled outputs, produced in response to the same question, as the baseline. More specifically, for each question $q$, GRPO samples a group of outputs $\{o_1, o_2, \cdots, o_G\}$  from the old policy  $\pi_{\theta_{old}}$  and then optimizes the policy model by maximizing the following objective:

\begin{equation}
\footnotesize
\begin{split}
    \mathcal{J}_{GRPO}(\theta) &= \mathbb{E}{[q \sim P(Q), \{o_i\}_{i=1}^G \sim \pi_{\theta_{old}}(O|q)]}  \\
    & \frac{1}{G}\sum_{i=1}^G \frac{1}{|o_i|} \sum_{t=1}^{|o_i|} \Bigg\{ \min \left[ \frac{\pi_\theta(o_{i,t} | q, o_{i,<t})}{\pi_{\theta_{old}}(o_{i,t} | q, o_{i,<t})} \hat{A}_{i,t},\, \text{clip} \left( \frac{\pi_\theta(o_{i,t} | q, o_{i,<t})}{\pi_{\theta_{old}}(o_{i,t} | q, o_{i,<t})},\, 1 - \varepsilon,\, 1 + \varepsilon \right) \hat{A}_{i,t} \right] \\
    &\hspace{8em} - \beta \mathbb{D}_{KL}\left[\pi_{\theta} || \pi_{ref}\right] \Bigg\},
\end{split}
\label{eq:GRPO-obj}
\end{equation}

where $\varepsilon$ and $\beta$ are hyper-parameters, and $\hat{A}_{i,t}$  is the advantage calculated based on relative rewards of the outputs inside each group only, which will be detailed in the following subsections. The group relative way that GRPO leverages to calculate the advantages, aligns well with the comparative nature of rewards models,  as reward models are typically trained on datasets of comparisons between outputs on the same question. Also note that, instead of adding KL penalty in the reward, GRPO regularizes by directly adding the KL divergence between the trained policy and the reference policy to the loss, avoiding complicating the calculation of  $\hat{A}_{i,t}$. And different from the KL penalty term used in (\ref{eq:PPO-reward}), we estimate the KL divergence with the following unbiased estimator \citep{kl_approx}:

\begin{equation}
\small
    \mathbb{D}_{KL}\left[\pi_{\theta} || \pi_{ref}\right] = \frac{\pi_{ref}(o_{i,t}|q,o_{i,<t})}{\pi_{\theta}(o_{i,t}|q,o_{i,<t})} - \log\left(\frac{\pi_{ref}(o_{i,t}|q,o_{i,<t})}{\pi_{\theta}(o_{i,t}|q,o_{i,<t})}\right) - 1,
\end{equation}

which is guaranteed to be positive.

\begin{algorithm}[t]
  \small
  \caption{Iterative Group Relative Policy Optimization}
  \textbf{Input} initial policy model $\pi_{\theta_{\text{init}}}$; reward models $r_\varphi$; task prompts $\mathcal{D}$; 
  hyperparameters $\varepsilon$, $\beta$, $\mu$
  \begin{algorithmic}[1]
    \State Initialize the policy model $\pi_\theta$ as $\pi_{\theta_{\text{init}}}$
    \For{iteration = 1, \dots, I}
       \State Set reference model $\pi_{ref}$ to $\pi_\theta$
      \For{step = 1, \dots, M}
      \State Draw a mini-batch $\mathcal{D}_b$ from $\mathcal{D}$
      \State Assign $\pi_{\theta_{old}} \gets \pi_\theta$
      \State For each query $q$ in $\mathcal{D}_b$, generate $G$ outputs $\{o_i\}_{i=1}^G \sim \pi_{\theta_{old}} (\cdot \mid q) $
      \State Evaluate each candidate output $o_i$ to obtain corresponding rewards $\{r_i\}_{i=1}^{G}$ using $r_{\varphi}$
      \State Estimate group-relative advantages $\hat{A}_{i,t}$ for every $t$-th token within each $o_i$
      \For{GRPO iteration = 1, \dots, $\mu$}
        \State Optimize the policy model $\pi_\theta$ according to the GRPO criterion (Equation \ref{eq:GC-GRPO})
      \EndFor
    \EndFor 
    \State Apply continual training to $r_\varphi$ utilizing a replay buffer.
    \EndFor 
  \end{algorithmic}
  \textbf{Output} $\pi_\theta$
  \label{alg:iter-grpo}
\end{algorithm}

\subsubsection{Outcome Supervision RL with GRPO} Specifically, for a given question $q$, a collection of responses $\{o_1, o_2, \cdots, o_G\}$ is generated utilizing the previous policy model $\pi_{\theta_{old}}$. These responses are subsequently evaluated by a reward model, which assigns a set of $G$ scores, denoted as $\mathbf{r} = \{r_1, r_2, \cdots, r_G\}$. The reward values are then standardized by removing the mean of the group and dividing by the group’s standard deviation. With outcome supervision, the normalized reward for each output $o_i$ is provided solely at the terminal step; the corresponding advantage values $\hat{A}_{i, t}$ for all tokens within $o_i$ are uniformly set to this normalized score, so that $\hat{A}_{i, t} = \widetilde{r}_i = \frac{r_i- {\rm mean}(\mathbf{r})}{{\rm std}(\mathbf{r})}$. Policy improvement is subsequently performed by maximizing the objective expressed in equation (\ref{eq:GRPO-obj}).

\subsubsection{Process Supervision RL with GRPO}  
While outcome supervision assigns a reward solely at the conclusion of each output, this approach often lacks the granularity needed to effectively guide policy learning in intricate mathematical problem solving. Motivated by \cite{wang2023math}, we investigate process supervision, in which feedback is provided at every reasoning stage. Specifically, for a given input question $q$ and a set of $G$ sampled responses $\{o_1, o_2, \cdots, o_G\}$, a process reward model evaluates each intermediate step, resulting in a set of per-step rewards: $\mathbf{R} = \{ \{r_1^{index(1)},\cdots,r_1^{index(K_1)}\}, \cdots,  \{r_G^{index(1)},\cdots,r_G^{index(K_G)}\} \}$. Here, $index(j)$ identifies the token position ending the $j$-th step, and $K_i$ denotes how many steps comprise the $i$-th response. To standardize the reward signals, we normalize them using the overall mean and standard deviation: $\widetilde{r}_i^{index(j)} = \frac{r_i^{index(j)} - {\rm mean(\mathbf{R})}}{{\rm std(\mathbf{R})}}$.

In this framework, the advantage attributed to each token is computed by aggregating the normalized rewards from subsequent steps: $\hat{A}_{i, t} = \sum_{index(j) \ge t} \widetilde{r}_i^{index(j)}$. The policy is then updated to maximize the objective specified in equation (\ref{eq:GRPO-obj}).

\subsubsection{Iterative RL with GRPO} During reinforcement learning, the previously trained reward model can become inadequate for guiding the evolving policy model. To address this, we investigate iterative RL with GRPO. As outlined in Algorithm \ref{alg:iter-grpo}, this approach regenerates reward model training datasets by sampling from the current policy model, and updates the reward model via continual training that maintains a 10\% subset of past data through a replay mechanism. Subsequently, the policy model is set as the new reference model and is further optimized using the updated reward model.

\subsection{Training and Evaluating DeepSeekMath-RL}

Our reinforcement learning experiments utilize the DeepSeekMath-Instruct 7B model. The RL training dataset is comprised of chain-of-thought-style problems derived from the GSM8K and MATH subsets present in the SFT data, totaling approximately 144,000 items. To assess how RL influences performance on benchmarks with limited or no exposure during RL training, we purposely omit additional SFT questions. The construction of the reward model’s training corpus adheres to the methodology outlined by \citep{wang2023math}. The initial reward model is initialized with DeepSeekMath-Base 7B, adopting a learning rate of $2 \times 10^{-5}$. When training with GRPO, the learning rate for the policy model is set to $1 \times 10^{-6}$, and we apply a KL divergence coefficient of 0.04. For each prompt, $64$ completions are generated. Both the maximum sequence length and the batch size for training are set at 1024. After every exploration phase, the policy model receives just a single update. For evaluation, DeepSeekMath-RL 7B is assessed on the same benchmarks used with DeepSeekMath-Instruct 7B. Within this framework, chain-of-thought GSM8K and MATH are treated as in-domain benchmarks for DeepSeekMath-RL 7B, while all remaining evaluations constitute out-of-domain assessment.

Table \ref{tab:sft_rl_math} presents a comparison of open- and closed-source models, evaluating their results with both chain-of-thought and tool-integrated reasoning approaches across English and Chinese benchmarks. Our observations are as follows: 1) Utilizing chain-of-thought reasoning, DeepSeekMath-RL 7B achieves accuracy rates of 88.2\% on GSM8K and 51.7\% on MATH. This places it ahead of all other open-source models ranging from 7B to 70B parameters and also surpasses most closed-source counterparts. 2) Notably, DeepSeekMath-RL 7B’s training is limited to chain-of-thought instruction tuning data from GSM8K and MATH, with its initialization from DeepSeekMath-Instruct 7B. Even within this narrow training regime, it consistently outperforms DeepSeekMath-Instruct 7B in all assessed metrics, underscoring the significant impact of reinforcement learning.

\section{Discussion}
This section presents the outcomes observed in our experiments involving pre-training and reinforcement learning (RL).

\subsection{Lessons Learnt in Pre-Training}

To begin, we discuss our observations from the pre-training phase. Throughout this section, except where noted, the training configurations described in Section \ref{sec:quality-policy} are employed. Additionally, references to the DeepSeekMath Corpus pertain to an 89B-token subset acquired during the second round of data gathering.

\subsubsection{Code Training Benefits Mathematical Reasoning}

An often-cited but unproven theory posits that practicing code enhances reasoning skills. In this work, we present a partial answer focused on mathematics: engaging in code training boosts the capacity of models to perform mathematical reasoning, whether or not external tools are employed.

To study how code training affects mathematical reasoning, we experimented with the following two-stage training and one-stage training settings:

\textbf{Two-Stage Training}

\begin{itemize}[topsep=0pt]
    \item \textbf{Code Training for 400B Tokens $\rightarrow$ Math Training for 150B Tokens}: Our approach first exposes DeepSeek-LLM 1.3B to 400B code tokens before proceeding with an additional 150B tokens specifically focused on mathematics;
    \item \textbf{General Training for 400B Tokens $\rightarrow$ Math Training for 150B Tokens}: For comparison, we conduct a control setup in which the initial 400B tokens come from general text (sampled from DeepSeek-AI’s extensive general corpus) rather than code, to examine whether training on code provides distinct benefits for mathematical reasoning performance relative to training on general content.
\end{itemize}

\textbf{One-Stage Training}

\begin{itemize}[topsep=0pt]
    \item \textbf{Math-Specific Training on 150B Tokens}: The DeepSeek-LLM 1.3B model is subjected to training on a dataset containing 150B tokens focused exclusively on mathematics;
    \item \textbf{Joint Training with 400B Code Tokens and 150B Math Tokens}: Post-code math finetuning has been observed to impair the model’s coding ability. To address this, we explore a unified training scheme in which code and math tokens (400B code, 150B math) are blended from the start, aiming both to boost mathematical reasoning and to mitigate the issue of catastrophic forgetting.
\end{itemize}

\paragraph{Results}
The downstream outcomes across various training configurations are presented in Table \ref{tab:code-math} and Table \ref{tab:code-math-general-eval}.

Incorporating code training enhances program-aided mathematical reasoning across both two-stage and one-stage training paradigms. Table \ref{tab:code-math} demonstrates that, in the context of two-stage training, code training by itself leads to considerable improvements in solving GSM8K and MATH tasks with Python. Introducing a subsequent math training phase brings about additional gains. Notably, in the one-stage training scenario, integrating code tokens with math tokens effectively addresses the problem of catastrophic forgetting observed in two-stage approaches, while also improving both coding abilities (refer to Table \ref{tab:code-math-general-eval}) and program-aided mathematical reasoning (see Table \ref{tab:code-math}).

Training on code similarly enhances mathematical reasoning in contexts where tools are not utilized. In the two-stage training paradigm, early exposure to code yields moderate gains even before focused mathematical instruction, and further facilitates the efficiency of later math-focused training, culminating in superior results. In contrast, when code and mathematics tokens are mixed during a single-stage training process, this seems to diminish the model’s mathematical reasoning capabilities without tool use. A possible explanation for this outcome is that DeepSeek-LLM 1.3B, with its relatively small parameter count, may be insufficiently robust to effectively learn from both code and mathematical data when presented together.

\subsubsection{ArXiv Papers Seem Ineffective in Improving Mathematical Reasoning}

ArXiv papers are commonly included as a component of math pre-training data \citep{minerva,gpt-f,llemma,mathpile}. However, detailed analysis regarding their impact on mathematical reasoning has not been extensively conducted. Perhaps counter-intuitively, according to our experiments, arXiv papers seem ineffective in improving mathematical reasoning. We experiment with models of different sizes, including DeepSeek-LLM 1.3B and DeepSeek-Coder-Base-v1.5 7B \citep{deepseek-coder}, using arXiv corpora that underwent varied processing pipelines:

\begin{itemize}[topsep=0pt]
    \item \textbf{MathPile} \citep{mathpile}: a dataset comprising 8.9 billion tokens, assembled using heuristic-based cleaning and filtering procedures, with scientific arXiv papers representing more than 85\% of the material;
    \item \textbf{ArXiv-RedPajama} \citep{redpajama}: a collection including all arXiv LaTeX sources, stripped of preambles, comments, macros, and bibliography sections, resulting in a total size of 28.0 billion tokens.
\end{itemize}

In our study, we conduct separate training of DeepSeek-LLM 1.3B using 150B tokens and DeepSeek-Coder-Base-v1.5 7B using 40B tokens on each arXiv dataset. The findings indicate that exposure to arXiv papers does not contribute significantly to enhancing mathematical reasoning skills. Models trained solely on the arXiv corpus either show negligible gains or, in some cases, perform worse across a range of mathematical evaluation benchmarks of varying complexity used in this analysis. The tested benchmarks encompass quantitative reasoning tasks such as GSM8K and MATH (see Table \ref{tab:arxiv-cot}), multiple-choice assessments like MMLU-STEM (see Table \ref{tab:arxiv-cot}), and formal mathematics problems as in miniF2F (see Table \ref{tab:arxiv_atp}).

However, this conclusion has its limitations and should be taken with a grain of salt. We have not yet studied:

\begin{itemize}[topsep=0pt]
    \item The influence of arXiv tokens on other mathematical tasks, not examined in this study, for instance, the task of informalizing theorems, which involves transforming formal statements or proofs into their informal counterparts;
    \item The interaction between arXiv tokens and additional forms of data;
    \item The extent to which the advantages offered by arXiv papers might emerge when training larger-scale models.
\end{itemize}

Thus, further exploration is required, which we leave for future studies.

\subsection{Insights of Reinforcement Learning}
\subsubsection{Towards to a Unified Paradigm} Here, we introduce a unified framework for examining various training approaches, including SFT, RFT, DPO, PPO, and GRPO, and undertake experiments to investigate the components of this unified framework. Broadly, the gradient with respect to the parameter $\theta$ for a given training strategy can be represented as:

\begin{equation}
    \nabla_{\theta}\mathcal{J}_{\textcolor{red}{\mathcal{A}}}(\theta) = \mathbb{E}[\underbrace{(q,o) \sim \textcolor{red}{\mathcal{D}}}_{Data \ Source}]\left( \frac{1}{|o|} \sum_{t=1}^{|o|}  \underbrace{GC_{{\mathcal{A}}}(q, o, t, \textcolor{red}{\pi_{{rf}}})}_{Gradient \ Coefficient}  \nabla_{\theta}\log \pi_{\theta}(o_t | q, o_{<t})\right).
\label{eq:objective}
\end{equation}

There exist three key components: 1) \textit{Data Source $\mathcal{D}$}, which determines the training data; 2) \textit{Reward Function $\pi_{{rf}}$}, which is the source of the training reward signal; 3) \textit{Algorithm $\mathcal{A}$}: which processes the training data and the reward signal to the gradient coefficient $GC$ that determines the magnitude of the penalty or reinforcement for the data. We analyze several representative methods based on such a unified paradigm:

\begin{itemize}
    \item \textbf{Supervised Fine-tuning (SFT)}: In SFT, a pretrained model is adapted using a dataset curated through human selection for supervised fine-tuning.
    \item \textbf{Rejection Sampling Fine-tuning (RFT)}: RFT involves an additional round of fine-tuning applied to the SFT model, employing responses produced by the SFT model in response to SFT prompts; only outputs passing correctness-based filtering are retained for further training.
    \item \textbf{Direct Preference Optimization (DPO)}: DPO advances the refinement of the SFT model by training on outputs sampled from the SFT model and optimizing using a pairwise DPO loss on augmented pairs.
    \item \textbf{Online Rejection Sampling Fine-tuning (Online RFT)}: In contrast to RFT, Online RFT starts from the SFT model and proceeds by incrementally fine-tuning the policy model with new samples generated online by the current policy itself.
    \item \textbf{PPO/GRPO}: PPO/GRPO commences training from the SFT model and further strengthens the policy model by reinforcement learning on outputs dynamically sampled from the ongoing policy.
\end{itemize}

Table \ref{tab:data_policy} provides an overview of the constituent elements of these approaches. For an in-depth explanation of the derivation procedures, consult Appendix \ref{app:analysis-rl}.

\paragraph{Observation about Data Source} We classify data sources into two types: online sampling and offline sampling. In the case of online sampling, training data is generated by the current policy model as it explores in real time. Conversely, offline sampling refers to data drawn from the outputs of the initial SFT model. Both RFT and DPO utilize an offline data approach, whereas Online RFT and GRPO are based on online data collection.

As depicted in Figure \ref{fig:rl-analysis}, our analysis reveals that Online RFT achieves markedly better performance than RFT across two benchmarks. In particular, although Online RFT and RFT demonstrate similar outcomes during the initial phase of training, Online RFT establishes a clear lead as training progresses, underscoring the effectiveness of the online approach. This phenomenon can be explained by the fact that, at the onset, both the actor and the SFT model are highly similar, leading to sampled data that differ only slightly. As training advances, the divergence in data generated by the actor becomes more pronounced, and sampling data in real time proves increasingly beneficial.

\paragraph{Observation about Gradient Coefficient} The algorithm utilizes the reward signal to compute the gradient coefficient for parameter updates. In our experimental setup, we decompose the reward function into two components: `Rule' and `Model'. The `Rule' component evaluates response quality by assessing answer correctness, while the `Model' component involves training a reward model that assigns scores to each response; the reward model is trained using labels generated from rule-based assessments. As presented in Equations \ref{eq:GC-RFT} and \ref{eq:GC-GRPO}, a central distinction between GRPO and Online RFT is that GRPO uniquely calibrates its gradient coefficient in accordance with the reward value produced by the reward model. This feature enables GRPO to provide varied reinforcement or penalization proportional to the magnitude of response quality. Conversely, Online RFT does not possess this adaptive mechanism: it applies uniform reinforcement to all responses deemed correct and does not penalize incorrect responses.

As illustrated in Figure \ref{fig:rl-analysis}, GRPO achieves better results than online RFT, underscoring the advantage of modifying positive and negative gradient coefficients. Moreover, GRPO+PS consistently outperforms GRPO+OS, which demonstrates the effectiveness of adopting step-aware, fine-grained gradient coefficients. Additionally, we investigate the impact of iterative RL; in our experimental setup, we apply two successive iterations. The results presented in Figure \ref{fig:iter-rl} reveal that iterative RL leads to considerable gains, most notably during the initial iteration.

\subsubsection{Why RL Works?} In this study, we apply reinforcement learning utilizing a selected portion of the instruction tuning dataset, resulting in marked improvements over the instruction-tuned baseline. To further elucidate the mechanism underlying this improvement, we analyze Pass@K and Maj@K accuracy metrics for both Instruct and RL-enhanced models across two evaluation benchmarks. According to the results depicted in Figure \ref{fig:maj-pass}, reinforcement learning consistently boosts Maj@K accuracy while leaving Pass@K relatively unchanged. This pattern suggests that reinforcement learning strengthens the resilience of the model's output distribution, primarily by increasing the likelihood that a correct answer appears within the TopK outputs, rather than by fundamentally advancing core model capabilities. In a similar vein, \citep{wang2023making} reported a \textbf{misalignment problem} affecting reasoning in SFT-trained models and demonstrated that reasoning performance could be improved through targeted preference alignment methods \citep{yuan2023rrhf,song2023preference,wang2023making}.

\subsubsection{How to Achieve More Effective RL?}
\label{sec:effec_rl}
Our results show that RL is highly effective for mathematical reasoning problems. Moreover, we introduce a comprehensive framework that encompasses various well-known training approaches, framing them as either direct or approximated RL methods. According to this framework, as outlined in Equation \ref{eq:objective}, there are three fundamental elements: Data Source, Algorithm, and Reward Function. We highlight possible avenues for further research involving these three aspects.

\paragraph{Data Source} The data source serves as the foundational input for all training methodologies. Within the realm of RL, we define the data source specifically as unlabeled questions together with their corresponding responses, which are generated by sampling from the policy model. For this study, we utilize only those questions originating from the instruction tuning phase, and employ a straightforward nucleus sampling approach to generate the responses. We hypothesize that this may partly explain why our RL pipeline achieves improvements solely in Maj@K performance. Going forward, we plan to extend our RL framework to evaluate its performance on question prompts drawn from out-of-distribution datasets, combined with \textbf{advanced sampling (decoding) strategies}, such as tree-search-based techniques \citep{yao2023tree}. Moreover, recent developments in \textbf{efficient inference techniques} \citep{xia-etal-2023-speculative,leviathan2023fast,kwon2023efficient,xia2024unlocking}, which are crucial for optimizing the exploration capabilities of policy models, are also recognized as playing a pivotal role in this context.

\paragraph{Algorithms} Algorithms utilize both input data and reward feedback to derive the appropriate gradient coefficients, subsequently adjusting the model parameters. According to Equation \ref{eq:objective}, contemporary approaches essentially place complete \textbf{TRUST} in the reward function’s signal when modifying the conditional probabilities associated with specific tokens. Nevertheless, such absolute reliance is problematic since the reward feedback cannot always be presumed dependable, particularly in the context of highly intricate tasks. For instance, even datasets like PRM800K \citep{lightman2023let}, which have undergone rigorous annotation by expert annotators, still exhibit about 20\% inaccurate labels\footnote{\url{https://github.com/openai/prm800k/issues/12\#issuecomment-1728491852}}. In light of this, we propose to investigate reinforcement learning algorithms designed for resilience to noisy or imperfect reward signals. We posit that alignment frameworks akin to \textbf{WEAK-TO-STRONG} \citep{burns2023weak} hold the potential to significantly redefine the nature of learning algorithms.

\paragraph{Reward Function} The reward function serves as the main source of training feedback. In reinforcement learning, this is most often represented by a neural reward model. We identify three key research directions for improving reward models: 1) \textbf{Enhancing the generalization capacity of the reward model.} For a reward model to be effective, it must generalize well to unseen questions and advanced model outputs; failure in this area risks RL simply reinforcing the existing distribution of large language models (LLMs) rather than advancing their overall performance; 2) \textbf{Capturing and representing the uncertainty within the reward model.} Accounting for uncertainty may provide critical insights for connecting the deficiencies of current reward models with the principles of weak-to-strong learning methods; 3) \textbf{Developing efficient strategies for constructing high-quality, process-level reward models} that are capable of delivering granular feedback throughout complex reasoning tasks \citep{lightman2023let,wang2023math}.

\section{Conclusion, Limitation, and Future Work}  
In this work, we introduce DeepSeekMath, a model that achieves state-of-the-art results among open-source models on the competitive MATH benchmark and closely rivals the performance of proprietary systems. Built upon DeepSeek-Coder-v1.5 7B, DeepSeekMath is continually trained on 500B tokens, of which a substantial portion—120B math-related tokens—are drawn from Common Crawl. Our thorough ablation analyses demonstrate that web-sourced pages are a highly valuable resource for mathematical data, whereas arXiv appears less advantageous than initially anticipated. We propose Group Relative Policy Optimization (GRPO), an adaptation of Proximal Policy Optimization (PPO), which enables substantial gains in mathematical reasoning proficiency while reducing memory usage. Experimental findings validate the effectiveness of GRPO, even as DeepSeekMath-Instruct 7B already attains impressive benchmark results. Additionally, we establish a unified framework for interpreting a range of methodologies and highlight multiple promising avenues for advancing reinforcement learning approaches.

While DeepSeekMath~demonstrates strong results on benchmarks requiring quantitative reasoning, its performance in geometry and theorem-proving falls short when compared to proprietary models. For example, our preliminary evaluation revealed the model’s inability to solve questions involving triangles and ellipses, possibly reflecting a data selection bias during pre-training and fine-tuning stages. Moreover, due to limitations in model size, DeepSeekMath~lags behind GPT-4 in few-shot learning; whereas GPT-4 benefits from few-shot prompting, DeepSeekMath~exhibits comparable outcomes across both zero-shot and few-shot settings. Moving forward, we intend to refine our data selection framework to assemble an even higher-quality pre-training dataset. Additionally, we plan to investigate new approaches (as outlined in Section \ref{sec:effec_rl}) to enhance the reinforcement learning processes for LLMs.

\end{document}